\begin{problem}[Warmup]{(a)}
    Can you trisect the angle $45^\circ$?
\end{problem}
\unsolved

\begin{problem}[Warmup]{(b)}
    Suppose you could square a particular circle. Could you then square any circle?
\end{problem}
\unsolved

\begin{problem}[Warmup]{(c)}
    Did we discuss a real number in this chapter that is \emph{not} constructible?
\end{problem}
\unsolved

\begin{problem}[Warmup]{(d)}
    Is $\sqrt{2 - \sqrt{5}}$ constructible? (Be careful!)
\end{problem}
\unsolved

\begin{problem}[Warmup]{(e)}
    Suppose that $\sqrt{\pi}$ were a constructible number. Why would this mean that we could square a circle of radius 1? (We'll see in Chapter 39 that $\sqrt{\pi}$ is not constructible.)
\end{problem}
\unsolved

\begin{problem}[Warmup]{(f)}
    Suppose that $\sqrt[3]{2}$ were a constructible number. Why would this mean that we could double a cube with edge length 1? (We'll see in Chapter 39 that $\sqrt[3]{2}$ is not constructible.)
\end{problem}
\unsolved

\begin{problem}[Warmup]{(g)}
    Euclid's first postulate states: ``Let the following be postulated: to draw a straight line from any point to any point.'' His second postulate says: ``To produce a finite straight line continuously in a straight line.'' With what tool do these postulates equip us? Why is the tool `idealized'?
\end{problem}
\unsolved

\begin{problem}[Warmup]{(h)}
    Euclid's third postulate states: ``To describe a circle with any center and any distance.'' With what tool does this postulate equip us? Why is the tool `idealized'?
\end{problem}
\unsolved

\begin{problem}{1}
    Give the construction for $a/b$ required in the Theorem 37.2.
\end{problem}
\unsolved

\begin{problem}{2}
    \begin{enumerate}
        \item[(a)] Construct $\sqrt{5}$, using the method described in the proof of the
        last theorem.
        \item[(b)] Now provide an easier construction of $\sqrt{5}$, by constructing the
        diagonal of an appropriately chosen rectangle.
    \end{enumerate}
\end{problem}
\unsolved

\begin{problem}{3}
    Explain the steps necessary to construct $\sqrt[4]{3} + 1$ and
    $\sqrt{\frac{5}{2} + \sqrt{7}}$, when starting with a line segment of length 1.
\end{problem}
\unsolved

\begin{problem}{4}
    Explain the steps necessary to construct $\sqrt{5 + 2\sqrt{6}}$ and
    $\sqrt{2} + \sqrt{3}$. Then show that these two numbers are the same.
\end{problem}
\unsolved

\begin{problem}{5}
    Show that $\sqrt{3 + 2\sqrt{2}}$ is of the form $a + b\sqrt{2}$, for some integers $a$
    and $b$. This exercise reinforces the lesson of Exercise 4: A given element of
    $\mathbb{K}$ may be constructible by quite a distinct list of steps!
\end{problem}
\unsolved

\begin{problem}{6}
    \begin{enumerate}
        \item[(a)] Perform the following construction with compass and straightedge,
        justifying each step: Take a line segment $\overline{AB}$, and construct on it a
        square $ABDC$. Then find the midpoint $E$ of $\overline{AC}$. Extend line segment
        $\overline{AC}$ so that $A$ is between $E$ and $H$, and
        $\overline{BE} = \overline{EH}$. (See the diagram in the book.)
        \item[(b)] Now apply the Pythagorean Theorem to the right triangle $EAB$ to show
        that
        \[\overline{AB}(\overline{AB} - \overline{AH}) = \overline{AH}^2.\]
        \item[(c)] If $|\overline{AB}| = 1$, what quadratic equation in
        $\rho = |\overline{AH}|$ do we obtain from part b?
        \item[(d)] Show that $1/\rho = \rho + 1$.
        \item[(e)] What is the value of the constructible number $\rho$? This number (or its
        reciprocal) is called the \textbf{Golden Section}, and the construction we have
        given for it appears as Proposition 11 in Book 2 of Euclid's \emph{Elements}.
    \end{enumerate}
\end{problem}
\unsolved

\begin{problem}{7}
    In this exercise you will show that the regular pentagon is constructible.
    \begin{enumerate}
        \item[(a)] In the diagram in the book (a regular pentagon with top vertex $A$ and
        bottom vertices $B$, $C$, with segments $AB$, $AC$, $BC$ drawn), show that
        $\angle BAC = 36^\circ$ and $\angle ABC = \angle ACB = 72^\circ$.
        \item[(b)] Let $E$ be the point where the bisector of $\angle ABC$ intersects the
        opposite side. Note that $\triangle ABC$ is similar to $\triangle BCE$. If
        $\overline{BC} = 1$ and $\overline{EC} = x$, show that $x^2 + x - 1 = 0$.
        \item[(c)] From Exercise 6c, you should have noted that the positive solution to
        this quadratic is the golden section, $\rho = \frac{1 + \sqrt{5}}{2}$. Now argue
        that the regular pentagon is constructible.
    \end{enumerate}
\end{problem}
\unsolved

\begin{problem}{8}
    Technically speaking, the compass given by Euclid's third postulate (see Exercise h)
    is \emph{collapsible}: it cannot be used directly to transfer distances from one line
    segment to another. Prove that distances can be transferred, with a collapsible compass
    and a straightedge, and so the modern compass is mathematically equivalent to the Greek
    collapsible compass.
\end{problem}
\unsolved
